\begin{claim}{D.2.12}{Proved}{The suspended stabilization unit has no retraction}
Let $B=K(\FF_2,2)$ and $A_r=Q\Sigma^rB$, where
$Q=\Omega^\infty\Sigma^\infty$. For every $r\geq1$ and every
pointed map
\[
 g:\Sigma A_r\longrightarrow\Sigma^{r+1}B,
\]
the homomorphism
\[
 \widetilde H_{r+3}(g;\FF_2)=0.
 \tag{D.2.12a}
\]
Consequently the suspended stabilization unit
\[
 \Sigma\eta:\Sigma^{r+1}B\longrightarrow\Sigma A_r
 \tag{D.2.12b}
\]
has no homotopy retraction. Its canonical stable retraction,
obtained from the suspension-spectrum/infinite-loop-space
counit, cannot be represented by a map of these spaces.
This conclusion does not assert that the entire mapping space
is contractible or determine $P_{A_r}\Sigma^{r+1}B$.
\end{claim}
\begin{proof}
All homology and cohomology in this proof have coefficients in
$\FF_2$. Put $X=\Sigma^rB$, $n=r+2\geq3$, and
$V=\widetilde H_*(X)$. Write $\operatorname{Sq}_*^i$ for the
homology operation dual to $\operatorname{Sq}^i$.

\medskip\noindent
\textbf{The four low-degree cohomology groups.}
Serre's calculation, including its explicit $K(\FF_2,2)$
example, gives
\[
 H^2(B)=\FF_2\{u\},\quad
 H^3(B)=\FF_2\{v\},\quad
 H^4(B)=\FF_2\{u^2\},\quad
 H^5(B)=\FF_2\{uv,w\},
 \tag{D.2.12c}
\]
where $v=\operatorname{Sq}^1u$ and
$w=\operatorname{Sq}^2v$; see
\cite[\S2, no.~8, Theorem~2 and no.~9, pp.~203--204]{I_Serre1953}.
Let $a,b,d\in V$ be the suspended homology classes dual to
$u,v,w$, respectively, with $d$ zero on the suspended class
$uv$. Their degrees are $n,n+1,n+3$, and
\[
 \operatorname{Sq}_*^1b=a,\qquad
 \operatorname{Sq}_*^2b=0,\qquad
 \operatorname{Sq}_*^1d=0,\qquad
 \operatorname{Sq}_*^2d=b.
 \tag{D.2.12d}
\]
For the third equality use
$\operatorname{Sq}^1(u^2)=0$; the other equalities follow
directly from \textup{(D.2.12c)}, naturality, and compatibility
of Steenrod operations with suspension. These compatibility and
Cartan formulas are recorded in
\cite[\S1, nos.~2.1--2.3, p.~199]{I_Serre1953}.
Moreover
\[
 \operatorname{Sq}_*^1c=0\quad(c\in V_{n+2}).
 \tag{D.2.12e}
\]
Indeed its value is a multiple of $b$, and applying
$\operatorname{Sq}_*^1$ again gives that multiple of $a$,
whereas $(\operatorname{Sq}_*^1)^2=0$.

\medskip\noindent
\textbf{The tensor algebra and the source classes.}
The James construction identifies
\[
 H_*(\Omega\Sigma X)\cong T(V)
       =\bigoplus_{k\geq0}V^{\otimes k}
 \tag{D.2.12f}
\]
as a Pontryagin algebra. The generators are the images of
$X\to\Omega\Sigma X$. This follows from the reduced-product
homology calculation and the canonical equivalence in
\cite[\S3, especially (3.3)--(3.4); Theorem~5.6,
Corollary~5.7 and Lemma~5.5]{I_James1955}.
Since $r\geq1$, the reduced diagonal of $X$ is zero in homology;
all of $V$ is therefore primitive in this Hopf algebra. The
diagonal on $T(V)$ is the multiplicative diagonal determined by
these primitives. Steenrod operations preserve tensor length:
they preserve $V$ by naturality of the one-letter inclusion,
and their effect on words is given by the Cartan formula.
Let $\ell_k:T(V)\to V^{\otimes k}$ be the projection onto
tensor length $k$.

Denote the images of $a,b,d$ under $X\to A_r$ by
$\bar a,\bar b,\bar d$. They are primitive and satisfy the
same Steenrod relations. The Pontryagin algebra $H_*(A_r)$
is commutative, since $A_r$ is an infinite loop space. Hence
the Cartan formula gives
\[
 \operatorname{Sq}_*^2(\bar b\bar d)=\bar b^2,
 \qquad
 \operatorname{Sq}_*^1(\bar b^2)
   =\bar a\bar b+\bar b\bar a=0.
 \tag{D.2.12g}
\]

Suppose that the adjunct $f:A_r\to\Omega\Sigma X$ of $g$
is nonzero on bottom homology. Then $f_*(\bar a)=a$ and
$f_*(\bar b)=b$: the groups in degrees $n$ and $n+1$ of
the target are one-dimensional, and
$\operatorname{Sq}_*^1b=a$. We will use only naturality of
the diagonal and Steenrod operations for $f$; multiplicativity
of $f_*$ is not required.

\medskip\noindent
\textbf{The coefficient forced by the diagonal.}
Let $\lambda:V\to\FF_2$ extract the coefficient of $b$,
and set $\mu=\lambda\operatorname{Sq}_*^2$.
Thus $\lambda$ is supported in degree $n+1$ and $\mu$ in
degree $n+3$. Relation \textup{(D.2.12e)} says
$\lambda\operatorname{Sq}_*^1=0$.
Put
\[
 t=\ell_1f_*(\bar d),\qquad
 z=\ell_2f_*(\bar b\bar d)\in V\otimes V.
\]
Naturality and preservation of tensor length imply
$\mu(t)=\lambda\ell_1f_*(\bar b)=1$.
Primitivity of $\bar b,\bar d$ gives
\[
 \overline\Delta(\bar b\bar d)
    =\bar b\otimes\bar d+\bar d\otimes\bar b.
\]
Projecting naturality of this equation onto tensor length
$(1,1)$ in the target gives
\[
 z+\tau z=b\otimes t+t\otimes b,
 \tag{D.2.12h}
\]
where $\tau$ interchanges the factors. Here only length-two
words contribute to the left side, because every letter is
primitive. Applying $\lambda\otimes\mu$ yields
\[
 (\lambda\otimes\mu+\mu\otimes\lambda)(z)=1.
 \tag{D.2.12i}
\]
On the other hand, the coefficient of the word $bb$ in
$\operatorname{Sq}_*^2z$ is exactly this expression: the
middle Cartan term vanishes because
$\lambda\operatorname{Sq}_*^1=0$. By
\textup{(D.2.12g)} and naturality, we have proved
\[
 \operatorname{coeff}_{bb}\bigl(\ell_2f_*(\bar b^2)\bigr)=1.
 \tag{D.2.12j}
\]
This argument also covers $n=3$. In that case $f_*(\bar d)$
may contain a length-two term $aa$; the projection defining
$t$ removes it, and it contributes nothing to
\textup{(D.2.12h)}.

\medskip\noindent
\textbf{The incompatible Bockstein.}
The length-two part of total degree $2n+2$ can contain only
$bb$ and terms $ac,ca$ with $c\in V_{n+2}$. Equations
\textup{(D.2.12d)--(D.2.12e)} give
\[
 \operatorname{Sq}_*^1(bb)=ab+ba,\qquad
 \operatorname{Sq}_*^1(ac)=
 \operatorname{Sq}_*^1(ca)=0.
\]
Thus \textup{(D.2.12j)} forces the length-two part of
$\operatorname{Sq}_*^1f_*(\bar b^2)$ to be the nonzero word
sum $ab+ba$. Other tensor lengths cannot cancel it. This
contradicts \textup{(D.2.12g)} and naturality, proving that
$f$ is zero on bottom homology.

The unit $X\to QX$ and the homology suspension
$H_n(\Omega\Sigma X)\to\widetilde H_{n+1}(\Sigma X)$
are isomorphisms in these bottom degrees, by connectivity and
the Hurewicz theorem. Adjunction therefore identifies the
bottom homology map of $f$ with that of $g$. This proves
\textup{(D.2.12a)}. The map \textup{(D.2.12b)} is a bottom
homology isomorphism, so it cannot have a homotopy retraction.
The stable counit is its stable retraction by the adjunction
triangle identity; any space-level representative would also
be a bottom homology isomorphism, which is impossible.

The proof above treats maps from one suspension of $A_r$.
The corollary below extends the low-degree homology conclusion
to every suspension, but does not control successive attachments
in nullification. No conclusion about the value of
$P_{A_r}\Sigma^{r+1}B$ is inferred.
\end{proof}

\paragraph*{Corollary to D.2.12: first cone attachments.}
\begingroup\itshape
Keep $r\geq1$, put $n=r+2$, and let $Y=\Sigma^{r+1}B$.
For every $k\geq0$ and every map $h:\Sigma^kA_r\to Y$,
\[
 \widetilde H_i(h;\FF_2)=0
 \qquad(n+1\leq i\leq n+3).
 \tag{D.2.12k}
\]
In degree $n+4$, its image is contained in the line dual to
the suspended class $uv\in H^5(B;\FF_2)$.
Thus $n+4=r+6$ is the first degree not excluded by these
computations; existence of a map hitting that line is not asserted.
If $C_h$ is the homotopy cofiber of $h$, then
\[
 \widetilde H_i(Y;\FF_2)\longrightarrow
 \widetilde H_i(C_h;\FF_2)
 \quad\hbox{is injective for }n+1\leq i\leq n+3.
 \tag{D.2.12l}
\]
The bottom integral homology homomorphism is also injective.
The same injections hold for the cofiber of any small wedge
of maps of the displayed form, including the first simultaneous
attachment of cones on all maps $\Sigma^kA_r\to Y$.
These statements concern the first attachment to $Y$; no
iteration through subsequent cone attachments is claimed.
\par\endgroup
\begin{proof}
First consider the bottom target degree $n+1$. For $k\geq2$
the source has no homology in that degree, and $k=1$ is
\textup{(D.2.12a)}. For $k=0$, restrict $h$ along the unit
$\eta:X=\Sigma^rB\to QX=A_r$. Since $X$ is
$(n-1)$-connected and $n\geq3$, the Freudenthal suspension
theorem \cite[Corollary~4.24]{FGHHatcher2002} makes this unit an isomorphism on homology in degree
$n+1$. If $h$ were nonzero in that degree, the restricted map
$j=h\eta:X\to Y$ would therefore satisfy
\[
 j^*(\sigma^{r+1}u)=\sigma^rv.
\]
Applying $\operatorname{Sq}^2$ gives
$j^*(\sigma^{r+1}u^2)=\sigma^rw$.
But $\operatorname{Sq}^1(u^2)=0$, whereas the Adem and
top-square identities give
\[
 \operatorname{Sq}^1w
 =\operatorname{Sq}^1\operatorname{Sq}^2v
 =\operatorname{Sq}^3v=v^2\ne0.
\]
Here the nonvanishing is part of Serre's polynomial-algebra
calculation, and the Adem identity is recorded in
\cite[\S4, no.~32.2, p.~224]{I_Serre1953}.
Naturality of $\operatorname{Sq}^1$ is a contradiction.

Write $a_Y,b_Y,c_Y$ for the classes in target degrees
$n+1,n+2,n+3$ dual to the suspended $u,v,u^2$.
They satisfy
\[
 \operatorname{Sq}_*^1b_Y=a_Y,
 \qquad\operatorname{Sq}_*^2c_Y=a_Y.
\]
The already proved vanishing in degree $n+1$ and naturality
force the image in each of the next two degrees to be zero.
In degree $n+4$ there are two basis classes, dual to $uv,w$.
The latter maps to $b_Y$ under $\operatorname{Sq}_*^2$ and
the former maps to zero. Hence the image can involve only
the $uv$-dual class. The homology exact sequence of the
cofiber proves \textup{(D.2.12l)}.

Finally $\widetilde H_{n+1}(Y;\mathbb Z)=\mathbb Z/2$.
The integral homology image of $h$ in that group is zero:
its reduction modulo $2$ injects into mod-$2$ homology by
the universal coefficient sequence, and
\textup{(D.2.12k)} annihilates it. The integral cofiber
sequence gives the asserted bottom injection. Homology takes
the indicated small wedge to a direct sum, so the same
vanishing and cofiber-sequence arguments prove the simultaneous
version. No integral
injection in degree $n+3$ is inferred from mod-$2$
injectivity, and the altered cofiber has not been shown to
satisfy the same Steenrod and loop-Hopf-algebra hypotheses
as $Y$.
\end{proof}
